\section{Damasus, Rome, and the Four Gospels}

The most securely bounded Damasian episode is Jerome's work on the Gospels in Rome between 382 and 384. The evidence is Jerome's preface beginning \emph{Novum opus}, addressed to Pope Damasus, together with his near-contemporary correspondence and the textual comparison of the resulting Latin tradition. The episode is important precisely because its limits can be stated: four Gospels, revised from existing Latin against Greek manuscripts, rather than a commissioned translation of the entire Bible (Jerome, \emph{Praefatio in Evangelia}; Letter 27; Houghton 2016, 31--39).

\subsection{A risky correction of familiar words}

Jerome represents Damasus as pressing him to adjudicate among differing Latin copies by recourse to Greek. He anticipates a predictable charge: readers who recognize a familiar Gospel but hear altered wording may call correction corruption. Letter 27 shows that this was not a merely imagined objection. Jerome's defense appeals to variant Latin manuscripts and to the Greek source.

The preface also states a restrained editorial policy. Jerome did not aim to replace every familiar Latin expression. He describes correction where a difference altered the sense while leaving much else in place. Modern comparison broadly supports an intervention that is more extensive in Matthew and decreases through Mark and Luke to John. The result is a revision whose continuity with Old Latin is as important as its correction from Greek (Houghton 2016, 31--39).

\evidenceframe{Jerome's \emph{Praefatio in Evangelia}, addressed to Damasus, and Letter 27 to Marcella.}{Jerome defended a revision of the four Gospels against older Greek manuscripts in 382--384 and expected resistance from readers of familiar Latin forms.}{The preface is a participant's programmatic defense, not a neutral commission register. It expressly bounds the promised work to the Gospels and cannot establish a whole-Bible mandate.}

\subsection{Text, order, and apparatus}

The revision was not only a sequence of altered words. Jerome placed the Gospels in the order Matthew, Mark, Luke, and John and transmitted Eusebian cross-reference apparatus that helped readers locate parallel passages. Such paratext matters historically: an edition can change how a text is studied even when its wording is conservatively revised. Later manuscripts could preserve, modify, or omit these features independently of their verbal text.

The Greek base cannot be identified with one surviving manuscript. Jerome speaks generally of older Greek copies, and the Latin outcome reflects an editorial act rather than a mechanical transcription. Textual criticism can compare his Latin with Greek traditions and infer relationships, but it does not recover a signed exemplar placed beside his desk. The proper historical claim is therefore comparative: his revision moved selected Latin readings toward Greek evidence available to him.

\subsection{What happened outside the Gospels}

Jerome's later Pauline commentaries sometimes criticize readings that became standard in the Vulgate Epistles and propose alternatives. This is difficult to reconcile with his authorship of their received revision. The Latin text of Acts, the Epistles, and Revelation also shows a technique and Greek basis distinct from the Gospels. Current reconstruction commonly places much of this work in a related late-fourth- or early-fifth-century revision, perhaps in Roman circles, but the reviser is unknown and more specific social attribution remains hypothetical (Houghton 2016, 40--43; Houghton 2023, introduction).

This negative result is substantive. An anonymous reviser is not a gap to be filled with Jerome because his name is available. It is evidence that the later Vulgate New Testament was already composite at its formation.

\claimcheck{``Damasus ordered the Vulgate Bible''}{The secure statement is narrower: Damasus sponsored or urged Jerome's revision of the four Gospels. Jerome's own preface names only those books. His later Old Testament translations arose in different circumstances, and the remaining New Testament cannot securely be assigned to him.}

\subsection{Early reception}

The revised Gospels gained users without immediately eliminating other texts. Augustine valued Jerome's work from Greek and used the revised Gospels in \emph{De consensu evangelistarum}. At the same time, the continued survival and correction of Old Latin Gospel books shows coexistence rather than instantaneous replacement. Liturgical hearing, private copying, commentary, and regional book production all had their own tempos.

The Damasian project nevertheless set a durable pattern. It joined respect for received Latin with appeal to source-language manuscripts, and it made textual correction an ecclesial as well as philological problem. Those same tensions intensified when Jerome moved from revising Latin by Greek to translating Old Testament books from Hebrew.
